\section{Resource Demo and Case Study}

The demo is designed as a diagnostic-resource check rather than a leaderboard.
It asks whether \tool can ingest real public \sid artifacts from both a
canonical residual-quantization line and a recent recommender-native line, then
produce comparable audit tables with explicit caveats. GRID provides the
open RQ-KMeans-style path~\cite{ju2025grid}; ReSID provides the GAOQ line and
processed Musical artifacts~\cite{liang2026resid}. Sanity tokenizers are kept
as controlled references for metric sensitivity.

\paragraph{Case-study protocol.}
The paper-facing case study uses the same 23,742 Musical items for both rows.
The GRID row is a controlled feature-text export: ReSID processed feature text
is embedded and passed through the official GRID-style residual k-means module.
The ReSID row is a bounded FAMAE-to-GAOQ export on the same item universe. This
protocol is intentionally modest. It gives reviewers a same-item diagnostic
contrast while avoiding the stronger, unsupported claim that the paper
reproduces every training choice of TIGER, GRID, or the full ReSID pipeline.

\begin{table}[t]
\caption{Same-item Musical artifact profiles on 23,742 items. Higher unique
full codes, D3 L1 co-occurrence prefix recall, and D4 tail unique-SID ratio are
better; lower D2 full-collision rate is better. D5a prefixes are level-wise
prefix counts, not latency. This is not a downstream leaderboard: GRID
feature-text uses ReSID processed feature text, and ReSID is a bounded GAOQ
export.}
\label{tab:musical-diagnostic}
\small
\setlength{\tabcolsep}{2.0pt}
\begin{tabular}{@{}lrrrrl@{}}
\toprule
System & Unique & D2 full & D3 L1 & D4 tail & D5a prefixes \\
\midrule
GRID feature-text & 3,749 & 0.9769 & 0.0552 & 0.3695 & 64/3440/3749 \\
ReSID bounded & 23,742 & 0.0000 & 0.1535 & 1.0000 & 32/1280/23742 \\
\bottomrule
\end{tabular}
\end{table}

Table~\ref{tab:musical-diagnostic} compares two exports on the same 23,742
Musical items. The GRID row is intentionally labeled as feature-text: it uses
processed feature text from the ReSID artifact path and the official
MiniBatchKMeans-style residual quantization module, not a faithful raw-text
TIGER reproduction. Under that controlled setup, the row exposes high collision
pressure and limited tail capacity. The ReSID/GAOQ row exports one full code
per item in this bounded run, yielding zero full-code collisions and higher
D3-L1 collaborative prefix recovery. The D5a prefix column reports level-wise
prefix counts, a structural proxy for trie expansion rather than real serving
latency.
The same GRID Musical feature-text export was rerun locally with seeds 43 and
44; across seeds 42--44, duplicate-\sid rate remains 0.8327--0.8421 and
full-collision rate remains 0.9751--0.9769, so the collision pressure in
Table~\ref{tab:musical-diagnostic} is not a single-seed artifact.

\paragraph{Diagnostic findings.}
The case study and controller tables give three concrete audit findings without
making a downstream leaderboard claim. First, collision pressure is a stable
artifact property in this controlled GRID row: most items share a complete
\sid{} even though the same interface can ingest a collision-free ReSID
mapping. A method-inspired collision controller adds a D2b check: collided
GRID pairs are 3.86x more likely to share a training user than
popularity-matched non-collision pairs, while a collision-heavy hash control
has only 1.19x lift. Thus raw collision volume and interaction-qualified
collision risk are not equivalent. Second, collision-free addresses and
collaborative-prefix alignment are different objectives. ReSID has no full
collisions and higher D3-L1 than the GRID row, while the repository sanity
table shows that a hand-designed category-prefix control can recover more
level-1 co-occurrence neighbors without being a learned recommender tokenizer.
Third, D4 and D5a keep capacity and structure separate: a capacity controller
at width 24 can preserve head uniqueness while leaving tail unique ratio at
0.028, and variable-depth controllers show that active prefix cost can differ
from a fixed maximum-depth schema. These are artifact mechanisms that aggregate
Recall@K or NDCG would not localize by themselves.

\paragraph{What the artifact tables add.}
The repository tables are not backup evidence for hidden claims; they provide
resource checks that do not fit the PDF. Sanity controls verify that D1--D5a
react to extreme tokenizers. GRID scale, DACT, and MovieLens tables document
larger-item handling, optional D6 churn, and non-Amazon portability. The
method-inspired controller tables test collision qualification, capacity
pressure, and variable-depth cost while staying separate from named-method
coverage. These tables support artifact usability; Table~\ref{tab:musical-diagnostic}
remains the only method-facing case study in the main paper.
